\section{Zadanie 11}
\autor{Mateusz Wojaczek}
Przeczytaj Sekcję 5.6.1 \textit{Random Graph Models} (modele losowych grafów,
podobieństwo do modelu kul i urn). Wypowiedz i udowodnij Lemat 5.14 dla
malejących własności grafów.

\subsection*{Rozwiązanie}
Graf jest \textbf{monotonicznie malejący},
dla dowolnych właściwości, gdy zachodzi dla
\[
	G = (V, E),
\]
to zachodzi też dla
\[
	G' = (V, E'),
\]
gdzie \(E' \subset E\).

\medskip

Niech:
\[
	P(n, N) \text{ -- prawdopodobieństwo, że własność zachodzi dla grafu } G_{n, N},
\]
\[
	P(n, p) \text{ -- prawdopodobieństwo, że własność zachodzi dla grafu } G_{n, p}.
\]

\medskip

Niech
\[
	p^{+} = \frac{(1+\epsilon)N}{\binom{n}{2}},
	\qquad
	p^{-} = \frac{(1-\epsilon)N}{\binom{n}{2}},
\]
dla \(0 < \epsilon < 1\).

\medskip

\textbf{Lemat.}
\text{Dla monotonicznie malejących grafów właściwości zachodzi:}
\[
	P(n, p^{+}) - e^{-O(N)} \le P(n, N) \le P(n, p^{-}) + e^{-O(N)}.
\]

\medskip

\textbf{(i) Dowód:} \(P(n, p^{+}) - e^{-O(N)} \le P(n, N)\)

\medskip

Niech \(X\) będzie zmienną losową oznaczającą liczbę krawędzi w grafie \(G_{n, p^{+}}\).

\[
	P(n, p^{+}) = \sum_{k=0}^{\binom{n}{2}} P(n, k) \cdot \Pr(X = k).
\]

Rozbijamy sumę:
\[
	P(n, p^{+}) = \sum_{k<N} P(n, k)\Pr(X=k) + \sum_{k \ge N} P(n, k)\Pr(X=k).
\]

Dla monotonicznie malejących własności mamy, że \(P(n, k) \le P(n, N)\) dla \(k \ge N\), zatem
\[
	P(n, p^{+}) \le \Pr(X < N) + P(n, N)\Pr(X \ge N) \le \Pr(X < N) + P(n, N).
\]

Ponadto:
\[
	\Pr(X < N) = \Pr\left(X < \frac{1}{1+\epsilon} \ev{X}\right)
	\le \Pr\left(X < (1 +\frac{\epsilon}{2}) \ev{X}\right)
	\le e^{-\frac{(1+\epsilon)\epsilon^2 N}{8}} = e^{-O(N)}.
\]

Zatem:
\[
	P(n, p^{+}) \le P(n, N) + e^{-O(N)},
\]
czyli
\[
	P(n, p^{+}) - e^{-O(N)} \le P(n, N).
\]
\[
	\boxed{\text{ckd}}
\]

\bigskip

\textbf{(ii) Dowód:} \(P(n, N) \le P(n, p^{-}) + e^{-O(N)}\)

\medskip

Niech \(X\) będzie zmienną losową liczby krawędzi w grafie \(G_{n, p^{-}}\).

\[
	P(n, p^{-}) = \sum_{k=0}^{\binom{n}{2}} P(n, k)\Pr(X=k)
	= \sum_{k \le N} P(n, k)\Pr(X=k) + \sum_{k > N} P(n, k)\Pr(X=k).
\]

Z monotoniczności mamy:
\[
	P(n, p^{-}) \ge P(n, N)\Pr(X \le N) \ge P(n, N) - \Pr(X > N).
\]

Z kolei:
\[
	\Pr(X > N) = \Pr\left(X > \frac{1}{1-\epsilon}\ev{X}\right)
	\le \Pr\left(X > (1+\epsilon)\ev{X}\right)
	\le e^{-\frac{(1-\epsilon)\epsilon^2 N}{3}} = e^{-O(N)}.
\]

Zatem:
\[
	P(n, p^{-}) \ge P(n, N) - e^{-O(N)},
\]
czyli
\[
	P(n, N) \le P(n, p^{-}) + e^{-O(N)}.
\]
\[
	\boxed{\text{ckd}}
\]